\section{Project Identity}
\label{sec:otel-weaver-exp-005-collector-to-intake-identity}

\subsection{Purpose}
\label{sec:otel-weaver-exp-005-collector-to-intake-identity-purpose}

EXP-005 defines and partially validates the OTel Collector ingestion boundary that routes live
telemetry feedstock into the process intelligence intake pipeline. Its purpose is to demonstrate
that raw OpenTelemetry spans emitted by arbitrary instrumented services can be gatekept,
normalized, and written as OCEL~2.0-structured JSON by a production-grade collector without
any modification to the upstream service. The experiment operationalises the doctrine that OTel
Weaver output is feedstock, not process evidence: feedstock is admitted or refused at the
collector boundary before it is ever allowed to enter the process intelligence conformance court.
The collector is therefore the first law-enforcing membrane in the full-lifecycle stack.

EXP-005 was listed as component~E in the \texttt{GGEN\_OTEL\_WEAVER\_PI\_ALIVE\_001.md}
checkpoint (SHA-256 sealed). Its PARTIAL status reflects a complete specification and a fully
manufactured harness, but the absence of a live Docker collector run and a materialised
\texttt{pi\_intake\_events.json} artifact.

\subsection{Repository Surface}
\label{sec:otel-weaver-exp-005-collector-to-intake-identity-repository}

All experiment artifacts reside under:
\begin{verbatim}
/Users/sac/process-intelligence/otel-weaver/experiments/exp-005-collector-to-pi-intake/
\end{verbatim}
Primary files are:
\begin{itemize}
  \item \texttt{README.md} --- full pipeline diagram, OCEL~2.0 JSON event schema, Docker launch
        command, and admissibility rules.
  \item \texttt{config/otel-collector-config.yaml} --- production-ready OTTL transform rules
        for the three-stage collector pipeline.
  \item \texttt{harness.py} --- synthetic load generator manufacturing three deliberately
        differentiated spans (Span~A, Span~B, Span~C) covering the full
        admission/rejection surface.
  \item \texttt{output/} --- intended destination of \texttt{pi\_intake\_events.json};
        currently empty because no live collector run has been executed.
\end{itemize}

Supporting doctrine and checkpoint files are found in:
\begin{verbatim}
/Users/sac/process-intelligence/otel-weaver/
  doctrine/otel-weaver-is-feedstock.md
  checkpoints/GGEN_OTEL_WEAVER_PI_ALIVE_001.md
  checkpoints/GGEN_OTEL_WEAVER_PI_RUNTIME_001.md
  checkpoints/GGEN_OTEL_WEAVER_PI_PARTIAL_001.md
  intel/collector-integration-model.yaml
  mappings/otel-to-pi-witness-map.yaml
  ggen/manifests/live-check-intake.manifest.yaml
\end{verbatim}

\subsection{Primary Research Role}
\label{sec:otel-weaver-exp-005-collector-to-intake-identity-role}

Within the dissertation, EXP-005 occupies the thesis role
\textit{live-telemetry-to-pi-intake-experiment}. It is the empirical counterpart to the
theoretical claim that a standards-compliant observability collector can serve as the
admission membrane for a process intelligence intake pipeline without requiring custom
middleware. EXP-005 is the first experiment in the OTel Weaver series to operate on live
OTLP traffic rather than static schema fixtures: EXP-001 through EXP-004 operate offline
(registry diff, residual typing, refusal manufacturing, witness lattice). EXP-005 is the
boundary where offline schema law meets online telemetry traffic.

The experiment contributes three specific research primitives: (1)~the filter/pi OTTL
processor pattern that enforces the \texttt{process.pi.instance\_id} namespace gatekeeper;
(2)~the transform/pi OTTL processor pattern that injects conformance metadata
(\texttt{pi.intake.ingested\_at}, \texttt{pi.intake.collector\_version}) without modifying
the upstream span; and (3)~the Span~A/B/C feedstock taxonomy that partitions live telemetry
into three legally distinct classes for downstream process mining.

\subsection{Key Artifacts}
\label{sec:otel-weaver-exp-005-collector-to-intake-identity-artifacts}

\begin{description}
  \item[\texttt{config/otel-collector-config.yaml}] The authoritative collector configuration
        specifying the OTLP gRPC (\texttt{:4317}) and HTTP (\texttt{:4318}) dual-protocol
        receiver endpoints, the \texttt{filter/pi} and \texttt{transform/pi} OTTL processors,
        and the \texttt{file/pi\_json} exporter writing to \texttt{output/pi\_intake\_events.json}.
  \item[\texttt{harness.py}] Python synthetic load generator that fires Span~A (fully
        compliant), Span~B (missing \texttt{process.pi.instance\_id}, dropped at filter),
        and Span~C (missing \texttt{process.pi.witness.hash}, admitted by collector but
        expected to fail the downstream live-check gate).
  \item[\texttt{output/pi\_intake\_events.json}] The OCEL~2.0 JSON intake artifact.
        PARTIAL: this file has not been materialised by a live run.
  \item[\texttt{live-check-intake.manifest.yaml}] ggen manifest referencing a BLAKE3 receipt
        store at \texttt{receipts/live-check-intake-receipt.json} that does not yet exist.
\end{description}
